\subsection{SUBALGEBRAS STABLE UNDER ad h}

Lemma 1. Let \(\mathrm{V}\) be a vector subspace of \(\mathfrak{g}\) and \(\mathrm{R(V)}\) the set of \(\alpha \in \mathrm{R}\) such that \(\mathfrak{g}^{\alpha} \subset \mathrm{V}\) . Then, \((\mathrm{V} \cap \mathfrak{h}) + \sum_{\alpha \in \mathrm{R(V)}} \mathfrak{g}^{\alpha}\) is the largest vector subspace of \(\mathrm{V}\) stable under \(\operatorname{ad} \mathfrak{h}\) .

A vector subspace W of V is stable under ad h if and only if

\[
W = (W \cap \mathfrak {h}) + \sum_ {\alpha \in R} (W \cap \mathfrak {g} ^ {\alpha})
\]

(Algebra, Chap. VII, §2, no. 2, Cor. 1 of Th. 1). The largest vector subspace of V stable under ad \(\mathfrak{h}\) is thus \((\mathrm{V} \cap \mathfrak{h}) + \sum_{\alpha \in \mathrm{R}} (\mathrm{V} \cap \mathfrak{g}^{\alpha})\) . But \(\mathrm{V} \cap \mathfrak{g}^{\alpha} = \mathfrak{g}^{\alpha}\) for \(\alpha \in \mathrm{R}(\mathrm{V})\) , and \(\mathrm{V} \cap \mathfrak{g}^{\alpha} = 0\) for \(\alpha \notin \mathrm{R}(\mathrm{V})\) since \(\dim \mathfrak{g}^{\alpha} = 1\) . Q.E.D.

For any subset P of R, put

\[
\mathfrak {g} ^ {\mathrm{P}} = \sum_ {\alpha \in \mathrm{P}} \mathfrak {g} ^ {\alpha} \quad \mathfrak {h} _ {\mathrm{P}} = \sum_ {\alpha \in \mathrm{P}} \mathfrak {h} _ {\alpha}.
\]

If \(\mathrm{P}\subset \mathrm{R}\) and \(\mathbf{Q}\subset \mathbb{R}\) , we clearly have

\[
\begin{array}{c} {[ \mathfrak {h}, \mathfrak {g} ^ {\mathrm{P}} ] \subset \mathfrak {g} ^ {\mathrm{P}}} \\ {[ \mathfrak {g} ^ {\mathrm{P}}, \mathfrak {g} ^ {\mathrm{Q}} ] = \mathfrak {g} ^ {(\mathrm{P+Q}) \cap \mathrm{R}} + \mathfrak {h} _ {\mathrm{P} \cap (- \mathrm{Q})}.} \end{array}\tag{1}
\]

\begin{enumerate}
\def\labelenumi{(\arabic{enumi})}
\setcounter{enumi}{1}
\tightlist
\item
\end{enumerate}

Recall (Chap. VI, §1, no. 7, Def. 4) that a subset P of R is said to be closed if the conditions \(\alpha \in P\) , \(\beta \in P\) , \(\alpha + \beta \in R\) imply \(\alpha + \beta \in P\) , in other words if \((P + P) \cap R \subset P\) .

Lemma 2. Let \(\mathfrak{h}'\) be a vector subspace of \(\mathfrak{h}\) and \(\mathrm{P}\) a subset of \(\mathrm{R}\) . Then \(\mathfrak{h}' + \mathfrak{g}^{\mathrm{P}}\) is a subalgebra of \(\mathfrak{g}\) if and only if \(\mathrm{P}\) is a closed subset of \(\mathrm{R}\) and

\[
\mathfrak {h} ^ {\prime} \supset \mathfrak {h} _ {\mathrm{P} \cap (- \mathrm{P})}.
\]

Indeed,

\[
[ \mathfrak {h} ^ {\prime} + \mathfrak {g} ^ {\mathrm{P}}, \mathfrak {h} ^ {\prime} + \mathfrak {g} ^ {\mathrm{P}} ] = [ \mathfrak {h} ^ {\prime}, \mathfrak {g} ^ {\mathrm{P}} ] + [ \mathfrak {g} ^ {\mathrm{P}}, \mathfrak {g} ^ {\mathrm{P}} ] = \mathfrak {h} _ {\mathrm{P} \cap (- \mathrm{P})} + [ \mathfrak {h} ^ {\prime}, \mathfrak {g} ^ {\mathrm{P}} ] + \mathfrak {g} ^ {(\mathrm{P} + \mathrm{P}) \cap \mathrm{R}}.
\]

Hence \(\mathfrak{h}' + \mathfrak{g}^{\mathrm{P}}\) is a subalgebra of \(\mathfrak{g}\) if and only if

\[
\mathfrak {h} _ {\mathrm{P} \cap (- \mathrm{P})} \subset \mathfrak {h} ^ {\prime} \text { and } \mathfrak {g} ^ {(\mathrm{P} + \mathrm{P}) \cap \mathrm{R}} \subset \mathfrak {g} ^ {\mathrm{P}}
\]

which proves the lemma.

PROPOSITION 1. (i) The subalgebras of g stable under ad h are the vector subspaces of the form \(h' + g^{P}\) , where P is a closed subset of R and \(h'\) is a vector subspace of h containing \(\mathfrak{h}_{\mathrm{P}\cap(-\mathrm{P})}\) .

\begin{enumerate}
\def\labelenumi{(\roman{enumi})}
\setcounter{enumi}{1}
\tightlist
\item
  Let \(\mathfrak{h}'\) , \(\mathfrak{h}''\) be vector subspaces of \(\mathfrak{h}\) and \(\mathrm{P},\mathrm{Q}\) closed subsets of \(\mathrm{R}\) , with \(\mathfrak{h}'\supset \mathfrak{h}_{\mathrm{P}\cap (-\mathrm{P})}\) , \(\mathfrak{h}''\subset \mathfrak{h}'\) and \(\mathrm{Q}\subset \mathrm{P}\) . Then \(\mathfrak{h}'' + \mathfrak{g}^{\mathrm{Q}}\) is an ideal of \(\mathfrak{h}' + \mathfrak{g}^{\mathrm{P}}\) if and only if
\end{enumerate}

\[
(P + Q) \cap R \subset Q \text { and } \mathfrak {h} _ {P \cap (- Q)} \subset \mathfrak {h} ^ {\prime \prime} \subset \bigcap_ {\alpha \in P, \alpha \notin Q} \operatorname{Ker} \alpha .
\]

Assertion (i) follows immediately from Lemmas 1 and 2. Let \(\mathfrak{h}',\mathfrak{h}''\) , P, Q be as in (ii). Then

\[
[ \mathfrak {h} ^ {\prime} + \mathfrak {g} ^ {\mathrm{P}}, \mathfrak {h} ^ {\prime \prime} + \mathfrak {g} ^ {\mathrm{Q}} ] = \mathfrak {h} _ {\mathrm{P} \cap (- \mathrm{Q})} + [ \mathfrak {h} ^ {\prime}, \mathfrak {g} ^ {\mathrm{Q}} ] + [ \mathfrak {h} ^ {\prime \prime}, \mathfrak {g} ^ {\mathrm{P}} ] + \mathfrak {g} ^ {(\mathrm{P} + \mathrm{Q}) \cap \mathrm{R}}.
\]

Hence, \(\mathfrak{h}'' + \mathfrak{g}^{\mathrm{Q}}\) is an ideal of \(\mathfrak{h}' + \mathfrak{g}^{\mathrm{P}}\) if and only if

\[
\mathfrak {h} _ {\mathrm{P} \cap (- \mathrm{Q})} \subset \mathfrak {h} ^ {\prime \prime}, [ \mathfrak {h} ^ {\prime \prime}, \mathfrak {g} ^ {\mathrm{P}} ] \subset \mathfrak {g} ^ {\mathrm{Q}}, \mathfrak {g} ^ {(\mathrm{P} + \mathrm{Q}) \cap \mathrm{R}} \subset \mathfrak {g} ^ {\mathrm{Q}}.
\]

This implies (ii).

PROPOSITION 2. Let \(\mathfrak{a}\) be a subalgebra of \(\mathfrak{g}\) stable under ad \(\mathfrak{h}\) , and let \(\mathfrak{h}' \subset \mathfrak{h}\) , \(\mathrm{P} \subset \mathrm{R}\) be such that \(\mathfrak{a} = \mathfrak{h}' + \mathfrak{g}^{\mathrm{P}}\) .

\begin{enumerate}
\def\labelenumi{(\roman{enumi})}
\item
  Let \(\mathfrak{k}\) be the set of \(x\in \mathfrak{h}'\) such that \(\alpha (x) = 0\) for all \(\alpha \in \mathrm{P}\cap (-\mathrm{P})\) . The radical of \(\mathfrak{a}\) is \(\mathfrak{k} + \mathfrak{g}^{\mathrm{Q}}\) , where \(\mathrm{Q}\) is the set of \(\alpha \in \mathrm{P}\) such that \(-\alpha \notin \mathrm{P}\) . Moreover, \(\mathfrak{g}^{\mathrm{Q}}\) is a nilpotent ideal of \(\mathfrak{a}\) .
\item
  a is semi-simple if and only if P = -P and \(h' = h_{P}\) .
\item
  \(\mathfrak{a}\) is solvable if and only if \(\mathrm{P} \cap (-\mathrm{P}) = \emptyset\) . In that case \([\mathfrak{a}, \mathfrak{a}] = \mathfrak{g}^{\mathrm{S}}\) , where
\end{enumerate}

\[
\mathrm{S} = ((\mathrm{P} + \mathrm{P}) \cap \mathrm{R}) \cup \{\alpha \in \mathrm{P} | \alpha (\mathfrak {h} ^ {\prime}) \neq 0 \}.
\]

\begin{enumerate}
\def\labelenumi{(\roman{enumi})}
\setcounter{enumi}{3}
\item
  \(\mathfrak{a}\) is reductive in \(\mathfrak{g}\) if and only if \(\mathrm{P} = -\mathrm{P}\) .
\item
  \(\mathfrak{a}\) consists of nilpotent elements if and only if \(\mathfrak{h}' = 0\) . Then \(\mathrm{P} \cap (-\mathrm{P}) = \emptyset\) , and \(\mathfrak{a}\) is nilpotent.
\end{enumerate}

We prove (v). If \(\mathfrak{a}\) consists of nilpotent elements, \(\mathfrak{a}\) is clearly nilpotent, and \(\mathfrak{h}' = 0\) since the elements of \(\mathfrak{h}\) are semi-simple. Assume that \(\mathfrak{h}' = 0\) . By Prop. 1 (i), \(P \cap (-P) = \varnothing\) . By Chap. VI, §1, no. 7, Prop. 22, there exists a chamber C of R such that \(P \subset R_{+}(C)\) . Hence, there exists an integer \(n > 0\) with the following properties: if \(\alpha_{1}, \ldots, \alpha_{n} \in P\) and \(\beta \in R \cup \{0\}\) , then

\[
\alpha_ {1} + \dots + \alpha_ {n} + \beta \notin \mathrm{R} \cup \{0 \}.
\]

This implies that every element of \(\mathfrak{g}^{\mathrm{P}}\) is nilpotent, hence (v).

We prove (iii). If \(\mathrm{P} \cap (-\mathrm{P}) = \varnothing\) , \(\mathfrak{g}^{\mathrm{P}}\) is a subalgebra of \(\mathfrak{g}\) (Prop. 1 (i)), and is nilpotent by (v). Now

\[
[ \mathfrak {a}, \mathfrak {a} ] = [ \mathfrak {h} ^ {\prime}, \mathfrak {g} ^ {\mathrm{P}} ] + [ \mathfrak {g} ^ {\mathrm{P}}, \mathfrak {g} ^ {\mathrm{P}} ] = [ \mathfrak {h} ^ {\prime}, \mathfrak {g} ^ {\mathrm{P}} ] + \mathfrak {g} ^ {(\mathrm{P} + \mathrm{P}) \cap \mathrm{R}} \subset \mathfrak {g} ^ {\mathrm{P}},
\]

so \(\mathfrak{a}\) is solvable and \([\mathfrak{a},\mathfrak{a}]\) is given by the formula in the proposition. If \(\mathrm{P}\cap (-\mathrm{P})\neq \emptyset\) , let \(\alpha \in \mathrm{P}\) be such that \(-\alpha \in \mathrm{P}\) . Then \(\mathfrak{h}_{\alpha} + \mathfrak{g}^{\alpha} + \mathfrak{g}^{-\alpha}\) is a simple subalgebra of \(\mathfrak{a}\) so \(\mathfrak{a}\) is not solvable.

We prove (i). Since P is closed, \((\mathrm{P} + \mathrm{Q}) \cap \mathrm{R} \subset \mathrm{P}\) . If \(\alpha \in P, \beta \in Q\) and \(\alpha + \beta \in R\) , we cannot have \(\alpha + \beta \in -P\) , for, P being closed, this would imply that \(-\beta = -(\alpha + \beta) + \alpha \in P\) whereas \(\beta \in Q\) ; thus, \((\mathrm{P} + \mathrm{Q}) \cap \mathrm{R} \subset \mathrm{Q}\) . This proves that \(g^{Q}\) is an ideal of a, nilpotent by (v). We have \(\mathrm{P} \cap (-\mathrm{Q}) = \varnothing\) , and \(\mathrm{P} \cap (-\mathrm{P}) = \mathrm{P} \cap \complement Q\) , so \(\mathfrak{h}_{\mathrm{P} \cap (-\mathrm{Q})} \subset \mathfrak{k} \subset \bigcap_{\alpha \in \mathrm{P}, \alpha \notin Q} \operatorname{Ker} \alpha\) . By Prop. 1 (ii),

\(\mathfrak{k} + \mathfrak{g}^{\mathrm{Q}}\) is an ideal of \(\mathfrak{a}\) . Since \(\mathrm{Q} \cap (-\mathrm{Q}) = \emptyset\) , this ideal is solvable by (iii). It is therefore contained in the radical \(\mathfrak{r}\) of \(\mathfrak{a}\) . Since \(\mathfrak{r}\) is stable under every derivation of \(\mathfrak{a}\) , \(\mathfrak{r}\) is stable under ad \(\mathfrak{h}\) . Hence there exists a subset \(S\) of \(P\) such that \(\mathfrak{r} = (\mathfrak{r} \cap \mathfrak{h}) + \mathfrak{g}^{\mathrm{S}}\) . Suppose that \(\alpha \in S\) and that \(-\alpha \in P\) . Then \(\mathfrak{h}_{\alpha} = [\mathfrak{g}^{\alpha}, \mathfrak{g}^{-\alpha}] \subset \mathfrak{r}\) , so \(\mathfrak{g}^{-\alpha} = [\mathfrak{h}_{\alpha}, \mathfrak{g}^{-\alpha}] \subset \mathfrak{r} = 0\) , so that \(-\alpha \in S\) ; by (iii), this contradicts the fact that \(\mathfrak{r}\) is solvable. Consequently, \(S \subset Q\) . Finally, if \(x \in \mathfrak{r} \cap \mathfrak{h}\) and if \(\alpha \in P \cap (-P)\) , then \([x, \mathfrak{g}^{\alpha}] \subset \mathfrak{g}^{\alpha} \cap \mathfrak{r} = 0\) , so \(\alpha(x) = 0\) ; this shows that \(x \in \mathfrak{k}\) . Hence \(\mathfrak{r} \subset \mathfrak{k} + \mathfrak{g}^{\mathrm{Q}}\) and the proof of (i) is complete.

We prove (iv). By (i), the adjoint representation of \(\mathfrak{a}\) on \(\mathfrak{g}\) is semi-simple if and only if \(\mathrm{ad}_{\mathfrak{g}}x\) is semi-simple for all \(x\in \mathfrak{k} + \mathfrak{g}^{\mathrm{Q}}\) (Chap. I, §6, no. 5, Th. 4); by (v), this is the case if and only if \(\mathrm{Q} = \varnothing\) , in other words \(\mathrm{P} = -\mathrm{P}\) .

We prove (ii). If \(\mathfrak{a}\) is semi-simple, \(P = -P\) by (i), so \(\mathfrak{h}_{\mathrm{P}} \subset \mathfrak{h}'\) ; further, \(\mathfrak{a} = [\mathfrak{a}, \mathfrak{a}] \subset \mathfrak{h}_{\mathrm{P}} + \mathfrak{g}^{\mathrm{P}}\) and consequently \(\mathfrak{h}' = \mathfrak{h}_{\mathrm{P}}\) . If \(P = -P\) and \(\mathfrak{h}' = \mathfrak{h}_{\mathrm{P}}\) , \(\mathfrak{a}\) is reductive by (iv), and \(\mathfrak{a} = \sum_{\alpha \in P} \mathfrak{s}_{\alpha}\) , so \(\mathfrak{a} = [\mathfrak{a}, \mathfrak{a}]\) and \(\mathfrak{a}\) is semi-simple.

PROPOSITION 3. Let \(\mathfrak{a}\) be a semi-simple subalgebra of \(\mathfrak{g}\) stable under \(\operatorname{ad}(\mathfrak{h})\) and let \(\mathrm{P}\) be the subset of \(\mathrm{R}\) such that \(\mathfrak{a} = \mathfrak{h}_{\mathrm{P}} + \mathfrak{g}^{\mathrm{P}}\) .

\begin{enumerate}
\def\labelenumi{(\roman{enumi})}
\item
  \(\mathfrak{h}_{\mathrm{P}}\) is a splitting Cartan subalgebra of \(\mathfrak{a}\) .
\item
  The root system of \((\mathfrak{a},\mathfrak{h}_{\mathrm{P}})\) is the set of restrictions to \(\mathfrak{h}_{\mathrm{P}}\) of elements of \(\mathrm{P}\) .
\end{enumerate}

Since \(h_{P}\) is stable under ad h, its normalizer in a is stable under ad h, and hence is of the form \(h_{P} + g^{Q}\) where \(Q \subset P\) (Lemma 1). If \(\alpha \in Q\) ,

\[
\mathfrak {g} ^ {\alpha} = \left[ \mathfrak {h} _ {\alpha}, \mathfrak {g} ^ {\alpha} \right] \subset \left[ \mathfrak {h} _ {\mathrm{P}}, \mathfrak {g} ^ {\alpha} \right] \subset \mathfrak {h} _ {\mathrm{P}},
\]

which is absurd. Thus \(Q = \varnothing\) and \(h_{P}\) is its own normalizer in a. This proves that \(h_{P}\) is a Cartan subalgebra of a. If \(x \in h_{P}\) , \(ad_{g}x\) , and a fortiori \(ad_{a}x\) , are triangularizable. Thus (i) is proved, and (ii) is clear.

By Prop. 1 (i), the subalgebras of g containing h are the sets \(h + g^{P}\) where P is a closed subset of R. By Chap. VII, §3, Prop. 3, every Cartan subalgebra of \(h + g^{P}\) is a Cartan subalgebra of g.

PROPOSITION 4. Let \(\mathfrak{a}\) be a subalgebra of \(\mathfrak{g}\) containing \(\mathfrak{h}\) , \(x\) an element of \(\mathfrak{a}\) , \(s\) and \(n\) its semi-simple and nilpotent components. Then \(s \in \mathfrak{a}\) and \(n \in \mathfrak{a}\) .

We have \((\operatorname{ad} x)\mathfrak{a} \subset \mathfrak{a}\) , so \((\operatorname{ad} s)\mathfrak{a} \subset \mathfrak{a}\) and \((\operatorname{ad} n)\mathfrak{a} \subset \mathfrak{a}\) . Since a is its own normalizer in g (Chap. VII, §2, no. 1, Cor. 4 of Prop. 4), \(s \in a\) and \(n \in a\) .

PROPOSITION 5. Let P be a closed subset of R.

\begin{enumerate}
\def\labelenumi{(\roman{enumi})}
\item
  \(\mathfrak{h} + \mathfrak{g}^{\mathrm{P}}\) is solvable if and only if \(\mathrm{P} \cap (-\mathrm{P}) = \varnothing\) . In that case, \([\mathfrak{h} + \mathfrak{g}^{\mathrm{P}}, \mathfrak{h} + \mathfrak{g}^{\mathrm{P}}] = \mathfrak{g}^{\mathrm{P}}\) .
\item
  \(\mathfrak{h} + \mathfrak{g}^{\mathrm{P}}\) is reductive if and only if \(\mathrm{P} = -\mathrm{P}\) .
\end{enumerate}

Assertion (i) follows from Prop. 2 (iii). If \(\mathrm{P} = -\mathrm{P}\) , \(\mathfrak{h} + \mathfrak{g}^{\mathrm{P}}\) is reductive (Prop. 2 (iv)). Assume that \(\mathfrak{a} = \mathfrak{h} + \mathfrak{g}^{\mathrm{P}}\) is reductive. Then

\[
\mathfrak {g} ^ {\mathrm{P}} = \left[ \mathfrak {h}, \mathfrak {g} ^ {\mathrm{P}} \right] \subset \left[ \mathfrak {a}, \mathfrak {a} \right] \subset \mathfrak {h} + \mathfrak {g} ^ {\mathrm{P}},
\]

so \([a, a]\) is of the form \(h' + g^{P}\) with \(h' \subset h\) ; since \([a, a]\) is semi-simple, P = -P (Prop. 2 (ii)).

